\documentclass[10pt,english]{article}

\usepackage{amsfonts,url,tikz,amsmath}

% Paper setup
\evensidemargin=0in
\oddsidemargin=0in
\textwidth=6.25in
\topmargin=-0.5in
\headheight=0.0in
\headsep=0.5in
\textheight=9.0in
\footskip=0.5in

\input{institution}

\newcommand{\Interior}[1]{\ensuremath{{#1}^{\circ}}}
\newcommand{\Closure}[1]{\ensuremath{\overline{#1}}}
\newcommand{\Complement}[1]{\ensuremath{{#1}^{c}}}

\newcommand{\Expect}{\ensuremath{\mathrm{E}}}
\newcommand{\vecnot}{\underline}
\newcommand{\RealNumbers}{\ensuremath{\mathbb{R}}}
\newcommand{\RationalNumbers}{\mathbb{Q}}
\newcommand{\ComplexNumbers}{\mathbb{C}}
\newcommand{\Real}{\mathrm{Re}}
\newcommand{\Span}{\mathrm{span}}
\newcommand{\Rank}{\mathrm{rank}}
\newcommand{\Nullity}{\mathrm{nullity}}
\newcommand{\Trace}{\mathrm{tr}}
\newcommand{\Diag}{\mathrm{diag}}
\newcommand{\dd}{\mathrm{d}}
\DeclareMathOperator*{\esssup}{ess\,sup}

% Use < , > inner product
\newcommand{\inner}[2]{{\left\langle #1 \mskip2mu , #2 \right\rangle}}
\newcommand{\tinner}[2]{{\langle #1 \mskip1mu , #2 \rangle}}

% Use < | > inner product
%\newcommand{\inner}[2]{{\left\langle #1 \mskip2mu \middle| \mskip2mu #2 \right\rangle}}
%\newcommand{\tinner}[2]{{\langle #1 \mskip1mu | \mskip1mu  #2 \rangle}}



\begin{document}

\title{ECE 586: Vector Space Methods \\ Lecture 22: Frequently Asked Questions}
\author{Henry D. Pfister \\ \univ}
\date{October 18, 2026}

\maketitle

\section{Alternating Projection}

\paragraph{What does alternating projection do for two closed subspaces?}
It projects onto each subspace in turn.
For closed subspaces $U,W$ of a Hilbert space, the sequence converges in norm to $P_{U\cap W}v_0$ (slide 1).
The intersection always contains zero, and each point in the intersection is fixed by both projections.

\paragraph{Why is a hyperplane usually not a subspace?}
For nonzero $w$, the set $\{v:\inner{v}{w}=0\}$ is a subspace.
Replacing zero by a nonzero constant $c$ shifts it away from the origin, producing an affine hyperplane.
Slides 2--3 derive its projection by shifting to a subspace, projecting, and shifting back.

\paragraph{How does the hyperplane projection formula work?}
For $H=\{u:\inner{u}{w}=c\}$ with $w\neq0$, it is
\[
P_H(v)=v-\frac{\inner{v}{w}-c}{\|w\|^2}w.
\]
The correction is parallel to the normal vector $w$, and substitution verifies that the new point satisfies the equality.
For example, projecting $(2,0)$ onto $x_1+x_2=1$ gives $(3/2,-1/2)$.
The course convention makes the inner product linear in its first argument.

\paragraph{Why can we solve a linear system by projecting onto individual equations?}
Each nonzero row of $Ax=b$ defines an affine hyperplane.
Their intersection is precisely the solution set.
Projecting cyclically onto these hyperplanes gives Kaczmarz's algorithm (slides 4--5).
The consistency assumption matters: an inconsistent system has no common point to which this feasibility argument can apply.

\paragraph{What is the Kaczmarz update, and what happens with a zero row?}
For a real row vector $a_i^T$, the relaxed update is
\[
v_{k+1}=v_k-s\frac{a_i^T v_k-b_i}{\|a_i\|_2^2}a_i,
\qquad i=(k\bmod m)+1.
\]
Here $s=1$ gives an exact projection and $0<s<1$ takes part of the correction.
Skip a zero row with $b_i=0$; a zero row with $b_i\neq0$ proves inconsistency.
The denominator requires a nonzero row.

\paragraph{Does a smaller relaxation parameter solve an inconsistent system in the least-squares sense?}
Not automatically. It can reduce oscillation, but a fixed relaxed cyclic projection scheme is not in general a least-squares solver.
For a least-squares or minimum-norm guarantee, use an algorithm designed for that objective, as discussed in Videos 18--20.
The suggestion on slide 5 that relaxation may help should not be read as such a guarantee.

\section{Convex Feasibility}

\paragraph{Do cyclic projections onto convex sets find the nearest point in their intersection?}
Generally no. For finitely many nonempty closed convex sets in a finite-dimensional space with nonempty common intersection, cyclic projections converge to a feasible point.
That point need not equal $P_{\cap_i C_i}v_0$.
This is the distinction emphasized on slide 6, ``Alternating Projection onto Convex Sets.''

\paragraph{What if the convex sets do not intersect, or the space is infinite dimensional?}
The finite-dimensional convergence statement on slide 6 assumes a nonempty intersection.
With an empty intersection, iterates can cycle or fail to converge to a feasible point.
In infinite-dimensional Hilbert spaces, the analogous convex-set results generally give weak convergence under the usual hypotheses, rather than the norm convergence guaranteed for two closed subspaces.

\paragraph{How do we project onto a half space?}
In a real Hilbert space, let $H=\{u:\inner{u}{w}\geq c\}$ with $w\neq0$.
If $v$ is feasible, leave it unchanged; otherwise project onto the boundary hyperplane:
\[
P_H(v)=v+\frac{\max\{0,c-\inner{v}{w}\}}{\|w\|^2}w.
\]
Slide 7 applies this idea to systems of linear inequalities.
For complex vectors an inequality must instead be specified using a real-valued expression, such as the real part of an inner product.

\end{document}
